\chapter{Heat-Kernel Area Term on a Causal Diamond}
\label{app:matrix_entropy}

This appendix records the standard heat-kernel contribution to the
entanglement entropy of a causal diamond. It does not derive the
equilibrium condition \(\delta S_D=\delta\langle H_D\rangle\) and does
not recover the Einstein equation from the matrix sector alone.

\section{Replica formula and cone expansion}
\label{sec:replica-cone}

The von Neumann entropy of the reduced state on a diamond \(D\) is
obtained from the replica trick,
\begin{equation}
  S_D=\bigl(1-n\partial_n\bigr)\ln Z(n)\Big|_{n=1},
\end{equation}
where \(Z(n)\) is the partition function on the \(n\)-sheeted cover
branched along the entangling surface \(\Sigma\). For a Laplace-type fluctuation operator on the emergent metric
(Assumption~\ref{ass:emergent-op}) the heat-kernel expansion on the
cone produces the area term
\begin{equation}
  S_D=\frac{\operatorname{Area}(\Sigma)}{4G_{\mathrm{eff}}}+S_{\mathrm{pot}},
\end{equation}
with the prefactor given by the Seeley--DeWitt coefficient that multiplies the
deficit. This is the standard cone computation for operators whose
endomorphism contains no explicit curvature coupling, on a replica geometry on
which the connection and the endomorphism stay bounded. For
\(E=E_0+QR\) with the conical delta of \(R\) not continued into \(E\), the
replica entropy differs from
\(\operatorname{Area}(\Sigma)/(4G_{\mathrm{eff}})\), with
\(G_{\mathrm{eff}}\) the curvature coefficient of Theorem~\ref{thm:Geff}, by
a contact term proportional to \(\operatorname{tr}Q\); if the delta is
continued into \(E\), the difference vanishes. For operators built from a
composite metric the curvature coefficient depends on the independent variables
chosen (Section~\ref{sec:status-B}). Chapter~\ref{ch:entanglement-equilibrium}
and Appendix~\ref{app:modular} use \(1/(4G_{\mathrm{eff}})\) as the name of
the coefficient of \(\delta A\). The conditional heat-kernel identity of Theorem~\ref{thm:Geff}
does not fix the numerical value of \(G_{\mathrm{eff}}\); the
coefficients \(A,B,C\) and physical matching remain open in
Problem~F.

\section{What is not derived}
\label{sec:not-derived}

\begin{itemize}
  \item The equilibrium condition
        \(\delta S_D=\delta\langle H_D\rangle\) under fixed volume.
        It is the hypothesis of Chapters~6--7 that forces \(m_T=0\).
        It is not obtained as a path-integral Ward identity from the
        replica variation alone.
  \item An evaluation of the fluctuation operator of the locked
        NAMCG condensate (including diffeomorphism and local-Lorentz
        ghosts) that would determine the precise endomorphism
        contribution to the area coefficient.
  \item Recovery of the infrared equation. Once the area law and the
        equilibrium condition are both assumed, the Einstein equation
        with a matrix source follows by the usual Jacobson argument;
        that assumption is not removed.
\end{itemize}

\section{Status}
\label{sec:appG-status}

The area term in \(S_D\) is the standard heat-kernel term for the operators of
Section~\ref{sec:replica-cone}. The
appendix adds no new derivation of the Ward identity, no scheme-fixed
coefficient, and no alteration to the open-problem ledger of
Appendix~\ref{app:status}.